\section{Supplementary Materials}
\subsection{Supplementary Methods}

\subsubsection{Deriving population-level quantities from individual-level components}

\subsubsection{Simulation approach}

We used stochastic, individual-based epidemic simulations to investigate the impact of individual-level infectiousness timing on epidemic dynamics. The simulation model has a modular design: an infection attempt generator produces a set of transmission attempts for each infected individual, and a population-level simulation script processes these attempts to propagate the epidemic. We implemented two simulation variants: a finite-population model with susceptible depletion, and an infinite-population model where every infection attempt succeeds. Both are implemented in {\bf \texttt{R}} (version 4.5.0), with performance-critical inner loops reimplemented in C++ via {\bf \texttt{Rcpp}} for speed.

\paragraph{Infection attempt generation.}
For each newly infected individual $i$ (with infection time $t_i$), the infection attempt generator proceeds as follows:
\begin{enumerate}
	\item \textbf{Draw the number of infection attempts.} The number of secondary infection attempts $\chi_i$ is drawn from a $\text{Poisson}(b_i)$ distribution, where $b_i$ is the individual's biological infectiousness. Throughout, we set $b_i \equiv R_0$ for all individuals, so that variation in realized transmission arises solely from infectiousness timing and contact dynamics.
	\item \textbf{Draw the latent period.} A latent period $l_i$ is drawn from a $\text{Gamma}((1-\psi)\alpha, \beta)$ distribution, representing the delay between infection and the onset of transmissibility. When $\psi = 1$, the latent period is zero; when $\psi = 0$, it accounts for the full variance of the generation interval.
	\item \textbf{Draw the timing of each infection attempt.} For each attempt $j = 1, \ldots, \chi_i$, an independent jitter $\epsilon_j$ is drawn from a $\text{Gamma}(\psi\alpha, \beta)$ distribution. The time of infection attempt $j$ relative to the index case's infection time is $\tau_{ij} = l_i + \epsilon_j$. When $\psi = 0$, $\epsilon_j = 0$ and all attempts occur simultaneously at time $l_i$; when $\psi = 1$, the $\epsilon_j$ are independent draws from the full $\text{Gamma}(\alpha, \beta)$ generation interval distribution.
	\item \textbf{Thin for time-varying contacts (if applicable).} When a time-varying contact process $c_i(t)$ is active, each infection attempt at calendar time $t_i + \tau_{ij}$ is accepted with probability $c_i(t_i + \tau_{ij}) / c_{\max}$, where $c_{\max}$ is the supremum of the contact process, using Poisson thinning. For the periodic contact model (Eq.~\ref{eq:periodic_contacts}), all individuals share the same deterministic contact function $c_i(t) = 1 - c_{\text{amp}} \cos(2\pi t / c_{\text{per}})$, and $c_{\max} = 1 + c_{\text{amp}}$. For the stochastic contact model, each individual's piecewise-constant contact trajectory is generated by drawing switching times from a Poisson process with rate $\lambda$ and contact levels from a $\text{Gamma}(\sigma, \sigma)$ distribution; the contact trajectory is generated over an interval $[0, 5\bar{g}]$ to ensure enough contact levels are drawn to accommodate all infection attempts. Thinning is performed against the realized maximum contact level over that individual's trajectory.
	\item \textbf{Thin for detect-and-isolate interventions (if applicable).} When a detect-and-isolate intervention is active, an individual-specific detection time $D_i$ is drawn (relative to peak infectiousness $l_i$) according to the detection mechanism in effect (symptom-based or screening-based). The individual participates with probability $\rho$. If participating, all infection attempts occurring after detection ($\tau_{ij} > t_i + D_i$) are removed with probability $\eta$, representing the effectiveness of post-detection isolation.
\end{enumerate}

\paragraph{Finite-population simulation.}
The finite-population simulation tracks susceptible depletion in a population of size $N$. It proceeds as follows:
\begin{enumerate}
	\item \textbf{Initialization.} A single index case is infected at time $t = 0$. All other $N-1$ individuals are susceptible.
	\item \textbf{Event scheduling.} Each infection attempt is stored as a pending event with its scheduled calendar time $t_i + \tau_{ij}$. Events are maintained in a priority queue, ensuring that the earliest pending event is always processed next.
	\item \textbf{Event processing.} At each step, the earliest event is removed from the queue. A target individual is selected uniformly at random from the population ($N - 1$ individuals excluding the infector). If the target is susceptible, infection occurs: the target is marked as infected, and its own infection attempts are generated and queued. If the target is already infected or recovered, the attempt fails (``wasted'' on a non-susceptible).
	\item \textbf{Termination.} The simulation terminates when the event queue is empty (epidemic extinction) or when a stopping criterion is met (\eg, a specified number of infections or calendar time is reached).
\end{enumerate}

\paragraph{Infinite-population simulation.}
The infinite-population simulation is a special case in which every infection attempt succeeds (no susceptible depletion). This is equivalent to a branching process approximation of the early epidemic and is useful for studying early-epidemic dynamics without the confounding effect of herd immunity. The implementation is identical to the finite-population version except that step 3 is simplified: every queued event produces a new infection, with no target selection or susceptibility check.

\paragraph{Default parameters.}
Unless otherwise stated, simulations used a population size of $N = 10{,}000$, with $n_{\text{sim}} = 5{,}000$ replicate simulations per scenario. Epidemic establishment was defined as reaching $0.05 N = 500$ cumulative infections. To cover a range of infectiousness profile shapes, we considered $\psi \in \{0, 0.2, 0.5, 0.8, 1\}$.

\paragraph{Benchmark pathogens.}
For illustration, we focused on three sets of epidemiological parameters reflecting published estimates of $R_0$ and $g(\tau)$ for influenza, SARS-CoV-2 (omicron), and measles. The generation interval distribution $g(\tau) \sim \text{Gamma}(\alpha, \beta)$ was parameterized using the mean generation time $\bar{g}$ and variance $\text{Var}[g]$ reported in the literature, with $\alpha = \bar{g}^2 / \text{Var}[g]$ and $\beta = \bar{g} / \text{Var}[g]$.

\begin{table}[h]
\centering
\begin{tabular}{lccccc}
\hline
Pathogen & $R_0$ & $\bar{g}$ (days) & $\text{Var}[g]$ (days$^2$) & $\alpha$ & $\beta$ \\
\hline
Influenza & 2 & 3.2 & 4.41 & 2.32 & 0.73 \\
SARS-CoV-2 (omicron) & 6 & 7.05 & 20.8 & 2.39 & 0.34 \\
Measles & 12 & 12.2 & 13.10 & 11.36 & 0.93 \\
\hline
\end{tabular}
\caption{Benchmark pathogen parameters used in simulations.}
\label{tab:benchmark_pathogens}
\end{table}




\clearpage


\subsection{Supplementary Figures}